


\section{问题重述}

\subsection{问题背景}


随着光伏等分布式能源快速发展，提高分布式能源本地消纳水平、增强系统灵活调节能力愈发重要。非孤岛式微网通过协调光伏发电、储能设备与外部电网购电，可在保障小区用电需求的同时促进本地消纳、减少并网压力并降低购电成本。


然而在实际运行中，光伏发电与小区负载存在供需错配，外部电网电价也会随时间波动，若调控不当将影响分布式能源消纳并增加购电成本。因此，如何合理协调光伏发电、储能充放电与外部电网购电，在保障供需平衡、促进分布式能源本地消纳和减少并网压力的同时，实现微网总购电费用最小化，具有重要的现实意义。



\subsection{问题要求}

本题层层递进，综合利用小区负载、光伏发电、储能状态及外部电网电价等信息，逐步考虑实际功率变化、预测更新与电价波动等因素，要求在满足小区用电需求的同时尽可能降低微网购电费用，制定合理的购电与储能充放电策略。

\textbf{问题1}
已知各时段外部电网电价、小区负载及全天光伏发电功率预测数据。在满足微网供电量不低于小区负载且储能设备日初与日末储电量相同的基础上，制定微网当天的计划购电策略，确定各时段购电量与储能充放电量，并计算全天购电量及购电费用。



\textbf{问题2}
考虑到小区负载和光伏发电功率随时间变化，且供电不足时需以5倍交易时刻电价紧急购电，制定每日计划购电策略，并确定指定日期的紧急购电时段及购电量。


\textbf{问题3}
引入每日四个时点更新的光伏发电功率预测，考虑因购电策略调整产生的违约及加购费用，制定相应的购电调整策略，并分析利用其他时刻预测信息进行策略调整的必要性。


\textbf{问题4}进一步考虑外部电网电价实时波动，在问题2、3的条件下重新制定微网购电策略并给出波动电价下相应的调控结果。

